\chapter{Giới thiệu đề tài}

\section{Đặt vấn đề}

Cuộc Cách mạng công nghiệp lần thứ tư cùng với sự phát triển nhanh chóng của trí tuệ nhân tạo (Artificial Intelligence - AI) đang tạo ra những thay đổi mạnh mẽ trong mọi lĩnh vực của đời sống kinh tế và xã hội. AI không còn chỉ là một công nghệ nghiên cứu mà đã trở thành công cụ quan trọng giúp doanh nghiệp nâng cao năng suất, tối ưu chi phí vận hành, tự động hóa quy trình nghiệp vụ và nâng cao khả năng cạnh tranh trên thị trường. Đặc biệt, sự xuất hiện của các mô hình AI tạo sinh (Generative AI) như ChatGPT, Gemini, Claude hay các mô hình mã nguồn mở đã giúp việc phát triển các giải pháp AI trở nên nhanh chóng, linh hoạt và dễ tiếp cận hơn bao giờ hết.
Tại Việt Nam cũng như trên thế giới, nhu cầu ứng dụng AI trong doanh nghiệp ngày càng gia tăng. Nhiều doanh nghiệp mong muốn triển khai các giải pháp như chatbot chăm sóc khách hàng, hệ thống phân tích dữ liệu, tự động hóa quy trình (AI Automation), xử lý tài liệu thông minh, trợ lý ảo, hệ thống đề xuất và các ứng dụng AI chuyên biệt khác nhằm nâng cao hiệu quả hoạt động và giảm chi phí nhân sự. Tuy nhiên, phần lớn doanh nghiệp, đặc biệt là các doanh nghiệp vừa và nhỏ, chưa sở hữu đội ngũ kỹ thuật có đủ chuyên môn để tự nghiên cứu và triển khai các hệ thống AI.
Bên cạnh đó, quá trình tìm kiếm đơn vị cung cấp dịch vụ AI hiện nay còn gặp nhiều khó khăn. Doanh nghiệp thường phải tìm kiếm thông qua các mối quan hệ cá nhân, mạng xã hội hoặc các nền tảng freelancer tổng quát. Những kênh này chưa được thiết kế riêng cho các dự án AI nên còn thiếu các cơ chế đánh giá chuyên môn, quản lý quy trình triển khai, theo dõi tiến độ và đảm bảo chất lượng dự án. Điều này làm tăng chi phí tìm kiếm đối tác, kéo dài thời gian triển khai và tiềm ẩn nhiều rủi ro trong quá trình hợp tác.
Ở chiều ngược lại, các chuyên gia AI, freelancer, startup công nghệ và các nhóm phát triển AI cũng gặp nhiều khó khăn trong việc tiếp cận khách hàng có nhu cầu thực sự. Mặc dù có năng lực chuyên môn cao, họ vẫn phải phụ thuộc vào các nền tảng việc làm chung hoặc hoạt động marketing cá nhân để tìm kiếm dự án. Điều này khiến việc xây dựng thương hiệu cá nhân, chứng minh năng lực và mở rộng thị trường trở nên khó khăn, đồng thời làm giảm cơ hội tiếp cận các dự án có giá trị lớn.
Ngoài ra, quá trình triển khai một dự án AI thường bao gồm nhiều giai đoạn như tiếp nhận yêu cầu, phân tích bài toán, báo giá, ký kết hợp đồng, trao đổi tài liệu, theo dõi tiến độ, nghiệm thu, thanh toán và đánh giá sau dự án. Hiện nay các bước này thường được quản lý thông qua nhiều công cụ riêng lẻ như email, Google Drive, Zalo, Slack hay các phần mềm quản lý công việc, dẫn đến dữ liệu bị phân tán, khó kiểm soát và làm giảm hiệu quả phối hợp giữa các bên.
Xuất phát từ những thực tiễn trên, nhóm nhận thấy cần có một nền tảng chuyên biệt dành riêng cho thị trường dịch vụ AI, đóng vai trò là cầu nối giữa doanh nghiệp và các chuyên gia AI. Nền tảng không chỉ giúp doanh nghiệp dễ dàng tìm kiếm đối tác phù hợp mà còn hỗ trợ các chuyên gia xây dựng hồ sơ năng lực, giới thiệu dịch vụ, nhận dự án và phát triển uy tín nghề nghiệp. Đồng thời, hệ thống cần tích hợp các chức năng quản lý toàn bộ vòng đời của dự án AI, bao gồm đăng yêu cầu, tìm kiếm chuyên gia, gửi đề xuất, trao đổi trực tuyến, quản lý tiến độ, thanh toán, đánh giá chất lượng và lưu trữ lịch sử hợp tác.
Từ những lý do trên, nhóm quyết định lựa chọn đề tài \textbf{AITasker – AI Marketplace Platform for AI Automation Services}. Hệ thống được định hướng xây dựng như một nền tảng Marketplace chuyên biệt dành cho lĩnh vực AI Automation, giúp kết nối hiệu quả giữa doanh nghiệp và các chuyên gia AI, nâng cao tính minh bạch trong quá trình hợp tác, tối ưu quy trình triển khai dự án và góp phần thúc đẩy quá trình chuyển đổi số trong doanh nghiệp. Đồng thời, đề tài cũng tạo điều kiện để nhóm vận dụng các kiến thức đã học về phát triển phần mềm, thiết kế hệ thống, cơ sở dữ liệu, điện toán đám mây và trí tuệ nhân tạo vào việc xây dựng một sản phẩm có tính ứng dụng thực tiễn cao.
\section{Lý do chọn đề tài}

Đề tài được lựa chọn dựa trên các lý do sau:

\begin{itemize}
    \item \textbf{Nhu cầu ứng dụng AI ngày càng gia tăng:} Trong bối cảnh chuyển đổi số diễn ra mạnh mẽ, trí tuệ nhân tạo (AI) đã trở thành một trong những công nghệ cốt lõi giúp doanh nghiệp nâng cao hiệu quả hoạt động. AI được ứng dụng trong nhiều lĩnh vực như chăm sóc khách hàng, phân tích dữ liệu, marketing, quản lý chuỗi cung ứng, sản xuất thông minh và tự động hóa quy trình nghiệp vụ. Việc ứng dụng AI không chỉ giúp tối ưu chi phí vận hành, giảm khối lượng công việc thủ công mà còn hỗ trợ doanh nghiệp đưa ra quyết định nhanh chóng và chính xác hơn, từ đó nâng cao năng lực cạnh tranh trên thị trường.
    \item \textbf{Thiếu nền tảng chuyên biệt cho dịch vụ AI:} Hiện nay, doanh nghiệp chủ yếu tìm kiếm chuyên gia AI thông qua các nền tảng freelancer tổng quát hoặc các kênh mạng xã hội. Tuy nhiên, những nền tảng này chưa được thiết kế riêng cho lĩnh vực AI nên thiếu các chức năng hỗ trợ đánh giá năng lực chuyên môn, quản lý dự án AI, báo giá và theo dõi tiến độ. Điều này khiến doanh nghiệp gặp nhiều khó khăn trong việc lựa chọn đối tác phù hợp, đồng thời các chuyên gia AI cũng khó tiếp cận đúng nhóm khách hàng có nhu cầu triển khai các giải pháp AI.
    \item \textbf{Xu hướng phát triển của AI và tự động hóa:} Sự phát triển nhanh chóng của các công nghệ như Machine Learning, Deep Learning, Generative AI, Large Language Models (LLMs) và AI Agents đang mở ra nhiều cơ hội để tự động hóa các quy trình nghiệp vụ trong doanh nghiệp. Các giải pháp AI Automation ngày càng được ứng dụng rộng rãi nhằm giảm chi phí, tăng năng suất và nâng cao chất lượng dịch vụ. Đây được xem là một trong những xu hướng công nghệ quan trọng, đóng vai trò thúc đẩy quá trình chuyển đổi số và phát triển kinh tế số trong tương lai.
    \item \textbf{Ý nghĩa thực tiễn của đề tài:} Việc xây dựng nền tảng \textbf{AITasker} góp phần giải quyết khoảng cách giữa doanh nghiệp có nhu cầu ứng dụng AI và các chuyên gia cung cấp dịch vụ AI. Hệ thống hỗ trợ doanh nghiệp dễ dàng đăng tải yêu cầu, tìm kiếm chuyên gia phù hợp, trao đổi thông tin, nhận báo giá, theo dõi tiến độ thực hiện và thanh toán trực tuyến. Đồng thời, nền tảng tạo môi trường minh bạch để các chuyên gia AI xây dựng hồ sơ năng lực, phát triển uy tín cá nhân và mở rộng cơ hội tiếp cận các dự án chất lượng.
    \item \textbf{Khả năng áp dụng kiến thức chuyên môn:} Đề tài là cơ hội để vận dụng tổng hợp các kiến thức đã học trong lĩnh vực Công nghệ thông tin như phân tích và thiết kế hệ thống, phát triển ứng dụng Web, lập trình Backend, thiết kế cơ sở dữ liệu, bảo mật hệ thống, điện toán đám mây, DevOps, API và tích hợp các dịch vụ AI hiện đại. Quá trình thực hiện đề tài giúp nhóm tích lũy kinh nghiệm xây dựng một hệ thống phần mềm hoàn chỉnh, có khả năng mở rộng, đáp ứng nhu cầu thực tế và phù hợp với xu hướng phát triển của ngành công nghệ.
\end{itemize}

\section{Mục tiêu đề tài}

Mục tiêu tổng quát của đề tài là xây dựng nền tảng \textbf{AITasker – AI Marketplace Platform for AI Automation Services}, đóng vai trò là cầu nối giữa doanh nghiệp có nhu cầu triển khai các giải pháp AI và các chuyên gia AI, freelancer hoặc nhóm phát triển. Hệ thống hướng đến việc hỗ trợ doanh nghiệp tìm kiếm chuyên gia phù hợp, quản lý toàn bộ quy trình thực hiện dự án và nâng cao hiệu quả hợp tác giữa các bên.
Để đạt được mục tiêu trên, đề tài đặt ra các mục tiêu cụ thể như sau:

\begin{itemize}
    \item Xây dựng nền tảng Marketplace chuyên biệt kết nối doanh nghiệp có nhu cầu triển khai các giải pháp AI với các chuyên gia AI, freelancer và đơn vị cung cấp dịch vụ AI, góp phần đơn giản hóa quá trình tìm kiếm và hợp tác giữa các bên.
    \item Xây dựng hệ thống quản lý tài khoản, xác thực người dùng và phân quyền theo từng vai trò (Doanh nghiệp, Chuyên gia AI và Quản trị viên), đảm bảo an toàn thông tin và kiểm soát quyền truy cập phù hợp với từng chức năng của hệ thống.
    \item Phát triển chức năng cho phép doanh nghiệp đăng tải yêu cầu hoặc dự án AI với đầy đủ thông tin như mô tả bài toán, mục tiêu, yêu cầu kỹ thuật, ngân sách, thời gian thực hiện và các tiêu chí lựa chọn chuyên gia.
    \item Xây dựng chức năng tìm kiếm, lọc và gợi ý các dự án phù hợp cho chuyên gia AI dựa trên lĩnh vực chuyên môn, kỹ năng, kinh nghiệm và lịch sử hoạt động, giúp nâng cao khả năng kết nối giữa doanh nghiệp và chuyên gia.
    \item Phát triển cơ chế gửi đề xuất thực hiện dự án (Proposal), báo giá, trao đổi thông tin và thương lượng giữa doanh nghiệp và chuyên gia trước khi ký kết hợp tác, tạo môi trường làm việc minh bạch và thuận tiện.
    \item Xây dựng hệ thống quản lý vòng đời dự án, hỗ trợ theo dõi tiến độ thực hiện, cập nhật trạng thái công việc, quản lý các mốc (Milestone), lưu trữ tài liệu và lịch sử trao đổi giữa các bên trong suốt quá trình triển khai.
    \item Xây dựng chức năng đánh giá, phản hồi và xếp hạng sau khi dự án hoàn thành nhằm tạo dựng uy tín cho doanh nghiệp và chuyên gia, đồng thời nâng cao chất lượng dịch vụ trên nền tảng.
    \item Nghiên cứu và tích hợp các kỹ thuật AI để hỗ trợ gợi ý chuyên gia phù hợp với yêu cầu dự án, đề xuất các dự án phù hợp với chuyên gia dựa trên hồ sơ năng lực, kỹ năng và lịch sử hợp tác, góp phần nâng cao hiệu quả kết nối.
    \item Thiết kế và triển khai hệ thống theo kiến trúc hiện đại, có khả năng mở rộng, đảm bảo tính bảo mật, hiệu năng, tính sẵn sàng và khả năng đáp ứng khi số lượng người dùng và dự án tăng lên trong tương lai.
\end{itemize}
\section{Đối tượng và phạm vi nghiên cứu}

\subsection{Đối tượng nghiên cứu}

Đề tài tập trung nghiên cứu các nội dung liên quan đến nền tảng Marketplace cung cấp dịch vụ AI, bao gồm mô hình kết nối giữa doanh nghiệp và chuyên gia AI, quy trình quản lý dự án, cơ chế trao đổi và cộng tác trực tuyến, cùng với các kỹ thuật gợi ý dựa trên AI để nâng cao hiệu quả kết nối.
Ngoài ra, đề tài cũng nghiên cứu các công nghệ phát triển ứng dụng Web hiện đại, hệ quản trị cơ sở dữ liệu, xác thực người dùng, bảo mật hệ thống và các giải pháp triển khai Backend phục vụ cho việc xây dựng một hệ thống hoàn chỉnh.

\subsection{Phạm vi nghiên cứu}

Đối tượng sử dụng hệ thống gồm:

\begin{itemize}
    \item \textbf{Doanh nghiệp:} Đăng tải dự án AI, tìm kiếm chuyên gia, trao đổi, quản lý tiến độ và đánh giá kết quả.
    \item \textbf{Chuyên gia AI/Freelancer:} Quản lý hồ sơ cá nhân, tìm kiếm dự án, gửi báo giá, trao đổi với khách hàng và thực hiện dự án.
    \item \textbf{Quản trị viên:} Quản lý người dùng, dự án, danh mục dịch vụ, xử lý báo cáo và giám sát hoạt động của hệ thống.
\end{itemize}

Trong phạm vi đề tài, hệ thống tập trung phát triển các chức năng chính sau:

\begin{itemize}
    \item Đăng ký, đăng nhập và xác thực người dùng.
    \item Quản lý hồ sơ cá nhân và hồ sơ năng lực.
    \item Đăng tải, chỉnh sửa và quản lý dự án AI.
    \item Tìm kiếm dự án và chuyên gia theo nhiều tiêu chí.
    \item Gửi, tiếp nhận và quản lý báo giá (Proposal).
    \item Chat và trao đổi trực tuyến giữa doanh nghiệp và chuyên gia.
    \item Theo dõi tiến độ thực hiện dự án.
    \item Đánh giá, nhận xét và xếp hạng sau khi hoàn thành dự án.
    \item Quản lý thông báo và lịch sử hoạt động.
    \item Chức năng quản trị hệ thống.
    \item Gợi ý chuyên gia hoặc dự án bằng AI.
\end{itemize}

Trong khuôn khổ đề tài, hệ thống tập trung xây dựng phiên bản Web phục vụ mục đích nghiên cứu và minh họa quy trình nghiệp vụ. Các chức năng thanh toán trực tuyến, ký hợp đồng điện tử hoặc triển khai ứng dụng di động sẽ được xem là hướng phát triển trong tương lai.

\section{Cấu trúc báo cáo}

Báo cáo được tổ chức thành sáu chương với nội dung chính như sau:

\begin{itemize}
    \item \textbf{Chương 1: Giới thiệu đề tài}
    Trình bày bối cảnh nghiên cứu, đặt vấn đề, lý do lựa chọn đề tài, mục tiêu nghiên cứu, đối tượng và phạm vi nghiên cứu, phương pháp thực hiện cũng như cấu trúc của báo cáo.

    \item \textbf{Chương 2: Khảo sát và phân tích yêu cầu hệ thống}
    Khảo sát thực trạng, phân tích bài toán, xác định các đối tượng sử dụng, xây dựng yêu cầu chức năng, yêu cầu phi chức năng và mô hình nghiệp vụ của hệ thống.

    \item \textbf{Chương 3: Phân tích và thiết kế hệ thống}
    Trình bày kiến trúc tổng thể của hệ thống, thiết kế cơ sở dữ liệu, xây dựng các sơ đồ UML như Use Case Diagram, Activity Diagram, Sequence Diagram, Class Diagram và thiết kế giao diện.

    \item \textbf{Chương 4: Xây dựng hệ thống}
    Giới thiệu các công nghệ được sử dụng, quá trình triển khai Backend, Frontend, cơ sở dữ liệu, API và xây dựng các chức năng chính của hệ thống.

    \item \textbf{Chương 5: Kiểm thử và đánh giá hệ thống}
    Thực hiện kiểm thử các chức năng của hệ thống, đánh giá kết quả đạt được, phân tích ưu điểm, hạn chế và mức độ đáp ứng các yêu cầu đã đặt ra.

    \item \textbf{Chương 6: Kết luận và hướng phát triển}
    Tổng kết những kết quả đạt được của đề tài, đánh giá đóng góp của hệ thống và đề xuất các hướng nghiên cứu, phát triển trong tương lai nhằm nâng cao tính ứng dụng của nền tảng.
\end{itemize}